% !Mode:: "TeX:UTF-8"
\fancyhead[CO]{\song\wuhao ~\chaptername \quad 实验结果~}
\fancypagestyle{plain}{\fancyhead[CO]{\song\wuhao ~\chaptername \quad 实验结果~}}

\chapter{实验结果}

% 第一节
\section{实验设置}
本章在 NVIDIA V100/A100 GPU 与海光 Z100/K100 DCU 异构平台上，对基于 DeepGen 生成的矩阵乘法算子、注意力机制及其对应的端到端模型进行基准性能测试。实验均在 Ubuntu 22.04 LTS 上进行，统一采用单精度浮点格式（FP32）作为计算精度，且未启用张量核硬件加速特性。在实验测试前均进行 50 次预热运行，随后执行 500 次测试，取平均值作为最终结果。本研究以 PyTorch 官方实现为参考基准计算相对 L2 范数误差，对矩阵乘法与注意力机制的输出进行数值精度验证。相关精度对比与分析结果已在本文第 \ref{precision} 节中详细说明。实验所用的硬件平台、软件依赖及环境配置详见表 \ref{tab:table6}。

\begin{table}[htbp]\wuhao
\caption{实验平台信息}\label{tab:table6}
\vspace{0.5em}
\renewcommand{\arraystretch}{1.35}
\begin{tabularx}{\textwidth}{ccccc}
\toprule[1.5pt]
平台 & \multicolumn{2}{c}{NVIDIA} & \multicolumn{2}{c}{海光} \\ \midrule[1pt]
型号 & V100-PCIe & A100-SXM & Z100 & K100 \\ 
显存 & 32G & 80G & 32G & 64G \\ 
频率 & 1380MHz & 1410MHz & 1319MHz & 2100MHz \\ 
编程模型 & CUDA 11.8 & CUDA 12.8 & ROCm 5.4、DTK 23.1 & ROCm 5.4、DTK 23.1 \\ 
依赖软件 & \multicolumn{2}{c}{\makecell[l]{PyTorch 2.6、TVM 0.15、Triton 2.3、\\ LLVM 21.0、IREE 3.9、Tilelang 0.1.6}} & \multicolumn{2}{c}{PyTorch 2.1、Triton 2.1、LLVM 21.0} \\ \bottomrule[1.5pt]
\end{tabularx}
\vspace{\baselineskip}
\end{table}

本文实验主要采用 FP32 精度对 DeepGen 生成算子的性能和数值正确性进行验证。选择 FP32 作为主要实验精度，一方面是为了在 NVIDIA GPU 和海光 DCU 平台上保持相对一致的实验条件，便于评估本文提出的双层 IR 表示、降级信息推断、跨平台优化 Pass 和自动调优框架本身的有效性；另一方面，FP16、BF16 等低精度高性能实现通常依赖目标硬件上的专用矩阵计算单元及其对应的指令映射机制，例如张量核或其他矩阵计算指令。该过程不仅涉及数据类型替换，还需要重新设计寄存器分块、共享内存布局、矩阵指令映射、累加精度选择和误差控制策略。因此，本文暂未将低精度混合精度路径作为主要实验内容，而是重点验证 DeepGen 在通用 FP32 代码生成路径上的正确性、性能收益和跨平台适配能力。

% =============================================================
%                4.2 矩阵乘法算子代码生成测试
% =============================================================
\section{矩阵乘法算子代码生成测试}

对于矩阵乘法的测试，本文首先将 DeepGen 生成的矩阵乘法算子与 NVIDIA 平台和海光平台上的 BLAS 库矩阵乘法算子进行对比，选取 $M$、$N$、$K$ 维度的 4 种不同比例进行测试；随后进一步对比不同平台上同样基于 MLIR 开发的代码生成工作 Triton 和经典代码生成开源项目 TVM，并在 $M$、$N$、$K$ 维度相等的方阵上进行实验。实验均采用每秒浮点计算次数 TFLOPS 作为性能指标，其中与 BLAS 库的对比实验还计算加速比作为衡量指标。

\begin{figure}[htbp]
\centering
\includegraphics[width=1\textwidth]{a100_matmul}
\caption{NVIDIA A100 上 DeepGen 与 cuBLAS 的矩阵乘法对比实验}\label{a100_matmul}
\vspace{\baselineskip}
\end{figure}

如图 \ref{a100_matmul} 所示，在 NVIDIA A100 平台上，DeepGen 在四组不同矩阵形状下整体均能达到与 cuBLAS 接近的性能水平，并在部分规模上实现超越。对于较小规模矩阵，由于计算量有限，线程块调度开销、访存启动开销以及边界处理代价占比较高，因此不同实现之间存在一定波动。随着矩阵规模增大，DeepGen 的性能逐渐稳定，并在多数测试点上与 cuBLAS 持平。尤其在部分大规模配置下，DeepGen 相比 cuBLAS 的性能提升超过 10\%，说明其自动生成代码在计算调度和数据访问组织方面已经具备较高质量。

分析可知，当 $M>10240$ 时，DeepGen 在 A100 上出现较明显的反超现象，这与 cuBLAS 在该规模区间内存在一定性能回落有关。其原因可能在于，cuBLAS 作为面向广泛场景的通用库，内部优化策略更强调整体适用性；而 DeepGen 通过块级分层 IR、降级信息推断和自动调优机制，能够生成更贴合当前矩阵形状与硬件结构的执行方案，因此在特定大规模场景下可以获得更优结果。

\begin{figure}[htbp]
\centering
\includegraphics[width=1\textwidth]{z100_matmul}
\caption{海光 Z100 上 DeepGen 与 rocBLAS 的矩阵乘法对比实验}\label{z100_matmul}
\vspace{\baselineskip}
\end{figure}

图 \ref{z100_matmul} 给出了海光 Z100 平台上的测试结果。总体来看，DeepGen 在该平台上同样表现出较强竞争力，在四组矩阵形状下均可达到与 rocBLAS 相近甚至更高的性能，部分测试点可达到 rocBLAS 的 103\%-124\%。这表明 DeepGen 的优化并非仅在 CUDA 平台上有效，而是在国产异构平台上同样具备较好的适配能力。

其原因在于，DeepGen 并不依赖特定厂商的专有模板，而是通过统一的块级 IR 显式描述矩阵乘法中的分块计算、线程映射和数据搬运关系，再借助跨平台优化 Pass 对共享存储访问、线程块计算顺序和访存离散化等细节进行调优。这类优化有助于减少 Bank 冲突并提升 Cache 命中率，因此能够在海光平台上取得优于官方库的结果。

需要说明的是，cuBLAS 和 rocBLAS 均为对应硬件平台上的厂商原生高性能库，其内部包含大量面向不同架构和矩阵形状的手工优化内核，通常可以代表标准矩阵乘法算子的强基线。因此，本文与厂商 BLAS 库的对比主要用于评估 DeepGen 自动生成代码相对于平台原生实现的竞争力，而不能简单理解为不同硬件平台之间的直接优劣比较。由于 NVIDIA CUDA 软件栈经过长期发展，在编译器、运行时、库函数启发式选择和性能分析工具等方面较为成熟；海光 DCU 相关的 ROCm/DTK 软件栈仍处于持续演进过程中，其库函数覆盖范围、形状启发式策略和后端优化成熟度可能与 CUDA 生态存在差异。因此，在部分规模下 DeepGen 超过 rocBLAS 或 cuBLAS，可能同时受到目标形状覆盖、库内部启发式选择和软件栈成熟度等因素影响。本文更关注的是 DeepGen 能否在不同厂商平台上通过统一 IR、可配置优化 Pass 和自动调优机制生成具有稳定竞争力的算子实现，从而降低跨平台迁移和手工优化成本。

在完成与官方 BLAS 库的比较后，本文进一步将 DeepGen 与 Triton、TVM 进行对比，结果如图 \ref{a100_ttd} 和图 \ref{z100_ttd} 所示。实验采用方阵 $M=N=K$ 配置，以减少形状差异带来的干扰。需要说明的是，部分规模下 TVM 未能搜索到可执行的矩阵乘法算子，因此图中存在空缺点。

\begin{figure}[htbp]
\centering
\includegraphics[width=0.6\textwidth]{a100_ttd}
\caption{NVIDIA A100 上 DeepGen 与 Triton 和 TVM 的矩阵乘法对比实验}\label{a100_ttd}
\vspace{\baselineskip}
\end{figure}

\begin{figure}[htbp]
\centering
\includegraphics[width=0.6\textwidth]{z100_ttd}
\caption{海光 Z100 上 DeepGen 与 Triton 和 TVM 的矩阵乘法对比实验}\label{z100_ttd}
\vspace{\baselineskip}
\end{figure}

从结果来看，DeepGen 在 NVIDIA 平台上的整体性能与 Triton、TVM 处于同一水平，并在部分规模下略有优势。在海光平台的小规模矩阵上，DeepGen 与 TVM、Triton 的性能处于同一量级；当矩阵规模增大后，DeepGen 相对于 Triton 和 TVM 的最高加速比分别达到 2.46 和 2.47 倍。这说明 DeepGen 在保持统一 IR 表达和跨平台生成能力的同时，仍能生成具有较强竞争力的高性能实现。总体而言，无论是与 cuBLAS、rocBLAS 的对比，还是与 Triton、TVM 的比较，DeepGen 在矩阵乘法代码生成上都取得了较理想的效果。该结果验证了前文所设计的分层 IR、优化 Pass 以及自动调优机制的有效性，也为后续注意力机制和端到端推理实验奠定了基础。
% =============================================================
%                4.3 多头注意力机制代码生成测试
% =============================================================
\section{多头注意力机制代码生成测试}
Transformer 架构已广泛应用于人工智能多个领域，其核心计算多头注意力机制的优化也是高性能推理中的重要问题。本节针对多头注意力机制进行测试对比，同时选取同样基于 MLIR 生态开发的端到端推理框架 IREE、以 TVM 为后端的 GPU 算子领域特定语言 Tilelang，以及 Triton 和 PyTorch 作为对比对象。由于海光平台暂未支持 Tilelang 生态，因此在海光 Z100 平台上仅与 IREE、Triton 和 PyTorch 进行对比。

在实验中，由于注意力批大小（Batch Size, BS）以及注意力头数（Head Number, HN）均由线程块的 $Z$ 和 $Y$ 维度进行索引，因此文中提及的批大小（BS）均为注意力批大小和注意力头数的乘积。为分析不同输入特征对性能的影响，本节分别开展了序列长度敏感性实验和批大小敏感性实验：前者固定批大小与头维度，考察序列长度变化对性能的影响；后者固定序列长度与头维度，分析批大小变化带来的吞吐率变化。

\begin{figure}[htbp]
\centering
\includegraphics[width=0.8\textwidth]{a100_attn_fix_bs}
\caption{NVIDIA A100 上多头注意力机制序列长度变化对比}\label{a100_attn_fix_bs}
\vspace{\baselineskip}
\end{figure}

\begin{figure}[htbp]
\centering
\includegraphics[width=0.8\textwidth]{z100_attn_fix_bs}
\caption{海光 Z100 上多头注意力机制序列长度变化对比}\label{z100_attn_fix_bs}
\vspace{\baselineskip}
\end{figure}
如图 \ref{a100_attn_fix_bs} 和图 \ref{z100_attn_fix_bs} 所示，本文首先测试了 A100 与 Z100 平台上多头对序列长度变化的敏感性。整体来看，DeepGen 在两个平台上的多数配置下都取得了最高或接近最高的 TFLOPS，尤其在小头维度场景中优势更为明显。在 A100 平台上，PyTorch 与 Tilelang 的整体性能较为接近，普遍优于 IREE 与 Triton；IREE 的性能随着注意力头维度（Head Dimension, HD）增大而有所提升，在较大头维度配置下逐渐超过 Triton，而 Triton 的性能曲线则相对平稳，没有随着参数规模增大而出现明显提升。而在 Z100 平台上，IREE 在小维度的较长注意力序列长度（Sequence Length, SL）的情况下，其性能有超越 Triton ，但是在大部分情况下其性能与其他框架相差较大；而 Triton 性能均处于稳定状态。

从加速比看，DeepGen 在 A100 上相对性能较优 PyTorch 框架，其平均加速在 $HD=64$ 时约为 1.47 倍，在 $HD=128$ 时约为 1.09 倍。这种差异与不同头维度下矩阵乘和 Softmax 计算的分块策略有关。对于 $HD=128$，现有库实现更容易获得较高效率；而在 $HD=64$ 时，资源更紧张，DeepGen 的融合实现能够更有效地减少中间结果写回并提升计算强度，因此优势更加明显；而相对于性能较差的 IREE 框架，在 $HD=64$ 的情况下，其平均加速比达到了 3.77 倍，但在 $HD=128$ 时，其加速比则为 1.77 倍，这是因为 IREE 在数据局部性利用与访存开销控制方面相对较弱，导致其难以充分发挥硬件计算能力，当 $HD$ 增大到 128，算子内部计算量增加，IREE 的硬件利用率有所提升，其固定编译与调度策略带来的劣势被部分摊薄，因此 DeepGen 相对其加速比有所下降。在 Z100 平台上，随着 $BS$ 增大，DeepGen 相对 PyTorch 的平均加速比为 1.22 倍，整体较为稳定，说明其优化策略在国产异构平台上同样具有较好的可迁移性；同时 DeepGen 相对于 IREE 的平均加速比高达 29.1 倍，这也印证了 IREE 依赖于固定平台，而无法做到有效的跨平台性能可移植，相反 DeepGen 在各平台上均能达到较高的性能。

\begin{figure}[htbp]
\centering
\includegraphics[width=0.8\textwidth]{a100_attn_fix_sl}
\caption{NVIDIA A100 上多头注意力机制批大小变化对比}\label{a100_attn_fix_sl}
\vspace{\baselineskip}
\end{figure}

\begin{figure}[htbp]
\centering
\includegraphics[width=0.8\textwidth]{z100_attn_fix_sl}
\caption{海光 Z100 上多头注意力机制批大小变化对比}\label{z100_attn_fix_sl}
\vspace{\baselineskip}
\end{figure}

在完成序列长度敏感性分析后，本文进一步开展了批大小敏感性实验，结果如图 \ref{a100_attn_fix_sl} 和图 \ref{z100_attn_fix_sl} 所示。总体来看，DeepGen 在 A100 和 Z100 平台上的大部分配置下仍然取得了最高 TFLOPS，且在长序列短头维度的场景下优势更加突出，尤其是在 $SL=16384$，$HD=64$ 的配置中表现最为明显。这说明随着批大小增加，注意力算子暴露出更高的块级并行度，而 DeepGen 设计的线程块映射与调度方式能够较充分地利用这部分并行性。

对于 A100 平台，IREE 在中大批量下表现出较好的扩展性和稳定性，但整体仍难以在多数配置下超过 DeepGen；Triton 的性能则相对平稳。DeepGen 相对 Triton 的平均加速约为 2.48 倍，相对 PyTorch 的平均加速约为 1.31 倍。对于 Z100 平台，当批大小较大时采用 $SL=512$ 配置，批大小较小时采用 $SL=16384$ 配置，其相对于 PyTorch 的加速比约为 1.12 倍至 1.38 倍；而在 $SL=512$ 和 1024 的配置下，其相对于 Triton 和 IREE 的加速比均在 2.81 倍以上，而在 $SL=16384$ 的长序列配置下，IREE 的性能提升明显，在 $BS=8$ 的情况下超过了 DeepGen，这说明 IREE 在长序列场景下对大规模并行计算与连续数据处理具有较好的适应能力。

综合图 \ref{a100_attn_fix_bs} 至图 \ref{z100_attn_fix_bs} 可以看出，DeepGen 在不同序列长度、不同头维度和不同批大小条件下都表现出较好的性能鲁棒性。其原因在于，DeepGen 通过分层 IR 对多头注意力机制的块级计算逻辑进行显式建模，并结合降级信息推断与优化 Pass，对线程束调度、共享存储访问和寄存器使用等细节进行针对性优化。相比之下，Triton 和 Tilelang 更依赖高级抽象的自动映射方式，在 CUDA 核心这类对细粒度调度高度敏感的场景中，其生成代码的性能容易受到限制；IREE 虽同样基于 MLIR，但采用更通用的编译路径，因此在小头维度场景下优势不够明显。

总体而言，本节实验从序列长度和批大小两个维度验证了 DeepGen 对注意力机制的代码生成能力。实验结果表明，DeepGen 在 NVIDIA A100 与海光 Z100 两个平台上均能在多数配置下取得最优或接近最优的性能，且在长序列、小头维度以及高并行度场景中优势更加突出。这说明本文提出的基于 MLIR 的注意力机制自动代码生成与优化方法是有效的，也进一步验证了前文分层 IR、优化 Pass 和自动调优策略在复杂 Transformer 核心算子中的适用性。

% --------------------------- Modify.3/7 ------------------------------
% =============================================================
%                     4.4 优化Pass消融实验
% =============================================================
\section{优化 Pass 消融实验}

为分析不同优化 Pass 对最终性能的影响，本文进一步设计了优化 Pass 开关消融实验。需要说明的是，本文提出的 Block Swizzle、Access Scatter、Software Pipeline 和 Local Split K 是作为可选优化策略被纳入自动调优空间。对于给定算子形状和目标硬件平台，DeepGen 可以选择启用或关闭其中部分 Pass。因此，本节将各优化 Pass 视为二值开关，通过比较不同 Pass 组合下的性能表现，分析各类优化策略的适用条件以及自动选择 Pass 组合的必要性。

本文在 NVIDIA A100 GPU 和海光 Z100 DCU 平台上，对多头注意力机制和矩阵乘法算子进行了优化 Pass 组合消融实验。其中，矩阵乘法算子的测试形状为 $M=N=K=2048$，多头注意力机制算子的测试形状为 $BHSD=(1,32,2048,64)$。算子生成过程中，除与消融实验相关的 Pass 配置外，其余调优参数均保持一致。实验结果如表 \ref{tab:matmul_pass_ablation} 和表 \ref{tab:mha_pass_ablation} 所示。其中，BS、AS、SP 和 LSK 的布尔取值分别表示是否启用 Block Swizzle、Access Scatter、Software Pipeline 和 Local Split K Pass 优化；全 0 或全 1 表示启用或不启用上述所有 Pass 优化；加粗数据表示由自动调优框架在候选空间中搜索得到的最优 Pass 组合。实验将全 0 配置作为基准，相对提升表示相对于基准的性能提升比例。最后，表中“--”表示该配置被省略、编译失败或执行错误。

\begin{table}[htbp]
\caption{矩阵乘法算子典型优化 Pass 组合消融实验}
\label{tab:matmul_pass_ablation}
\vspace{0.5em}\centering\wuhao
% \renewcommand{\arraystretch}{1.1}
\begin{tabular}{ccccccc}
\toprule[1.5pt]
平台 & BS & AS & SP & LSK & 性能（TFLOPS） & 相对提升 \\
\midrule[1pt]
\multirow{7}{*}{NVIDIA A100}
& 0 & 0 & 0 & 0 & 11.98 & -- \\
& 1 & 0 & 0 & 0 & 13.28 & 10.86\% \\
& 0 & 1 & 0 & 0 & 12.02 & 0.37\% \\
& 0 & 0 & 1 & 0 & 12.07 & 0.82\% \\
& 0 & 0 & 0 & 1 & 12.39 & 3.48\% \\
& 1 & 1 & 1 & 1 & 15.87 & 32.51\% \\
& \textbf{1} & \textbf{1} & \textbf{1} & \textbf{1} & \textbf{15.87} & \textbf{32.51\%} \\
\multirow{7}{*}{海光 Z100}
& 0 & 0 & 0 & 0 & 7.05 & -- \\
& 1 & 0 & 0 & 0 & 7.30 & 3.47\% \\
& 0 & 1 & 0 & 0 & 7.07 & 0.29\% \\
& 0 & 0 & 1 & 0 & 6.63 & -6.01\% \\
& 0 & 0 & 0 & 1 & -- & -- \\
& 1 & 1 & 1 & 1 & -- & -- \\
& \textbf{1} & \textbf{0} & \textbf{1} & \textbf{0} & \textbf{8.71} & \textbf{23.45\%} \\
\bottomrule[1.5pt]
\end{tabular}
\vspace{\baselineskip}
\end{table}

从实验结果可以看出，不同优化 Pass 组合对性能的影响并不具有单调叠加特征。表 \ref{tab:matmul_pass_ablation} 展示了矩阵乘法算子的消融实验结果。Block Swizzle 通常能够在较低资源开销下改善线程块访问顺序，从而提升 Cache 命中率，因此该 Pass 在大多数情况下都具有稳定的性能收益，在 NVIDIA 和海光平台上分别取得了 10.86\% 和 3.47\% 的性能提升。相比之下，Access Scatter 会修改线程的访存模式，当离散化宽度能够满足最大向量化访存宽度且不会产生 Bank 冲突时，能够带来一定性能收益；但如果离散化宽度不足，则可能增加访存指令数量，导致性能收益有限甚至出现性能下降。因此，在当前测试配置下，Access Scatter Pass 在两个平台上的性能提升并不明显。

此外，Software Pipeline 和 Local Split K 虽然能够提高访存隐藏能力和规约轴并行度，但也会引入额外的共享内存占用、同步开销以及局部规约开销。在 NVIDIA 平台上，由于硬件资源较为充足，上述 Pass 仍能够带来一定性能提升；而在海光平台上，由于共享内存容量和线程驻留资源更加敏感，启用相关 Pass 后可能导致共享内存超限、占用率降低甚至编译失败。因此，在海光平台上，自动调优器最终仅选择启用 Block Swizzle 和 Software Pipeline Pass，从而获得最优性能。

\begin{table}[htbp]
\caption{多头注意力机制算子典型优化 Pass 组合消融实验}
\label{tab:mha_pass_ablation}
\vspace{0.5em}\centering\wuhao
% \renewcommand{\arraystretch}{1.1}
\begin{tabular}{cccccc}
\toprule[1.5pt]
平台 & AS & SP & LSK & 性能（TFLOPS） & 相对提升 \\
\midrule[1pt]
\multirow{7}{*}{NVIDIA A100}
& 0 & 0 & 0 & 8.72 & -- \\
& 1 & 0 & 0 & 8.52 & -2.24\% \\
& 0 & 1 & 0 & 11.21 & 28.53\% \\
& 0 & 0 & 1 & 9.05 & 3.79\% \\
& 1 & 1 & 1 & -- & -- \\
& \textbf{0} & \textbf{1} & \textbf{0} & \textbf{11.21} & \textbf{28.53\%} \\
\multirow{7}{*}{海光 Z100}
& 0 & 0 & 0 & 3.56 & -- \\
& 1 & 0 & 0 & 3.41 & -4.21\% \\
& 0 & 1 & 0 & 4.21 & 18.26\% \\
& 0 & 0 & 1 & 3.68 & 3.37\% \\
& 1 & 1 & 1 & -- & -- \\
& \textbf{0} & \textbf{1} & \textbf{0} & \textbf{4.21} & \textbf{18.26\%} \\
\bottomrule[1.5pt]
\end{tabular}
\vspace{\baselineskip}
\end{table}

表 \ref{tab:mha_pass_ablation} 展示了多头注意力机制算子的消融实验结果。由于该算子的线程块在计算过程中仅沿单一维度进行映射，因此无法执行 Block Swizzle Pass 优化。从实验结果可以看出，Access Scatter Pass 在两个平台上均表现出负收益。这是因为当前配置下的访存模式已经能够较好地满足向量化访存需求，而 Access Scatter 会改变线程访问模式，导致离散化宽度无法满足最大向量化访存宽度，从而增加访存指令数量和访存开销，最终导致性能下降。

相比之下，Software Pipeline Pass 在两个平台上均取得了最明显的性能提升，在 NVIDIA A100 和海光 Z100 平台上分别获得了 28.53\% 和 18.26\% 的性能提升。这是因为多头注意力机制算子包含 $QK$ 矩阵乘、Softmax 归一化以及 $PV$ 矩阵乘等多个阶段，其性能瓶颈更多来源于全局访存延迟和多阶段流水之间的同步等待。Software Pipeline 能够通过重叠数据加载与计算过程，有效隐藏访存延迟，因此在多头注意力机制中具有更明显的优化效果。而 Local Split K Pass 仅能够带来较小幅度的性能提升，原因在于多头注意力机制的规约维度较小，其并行规约收益有限，同时局部规约过程还会引入额外同步与共享内存开销。最后，在两个平台上，同时启用所有 Pass 后均出现共享内存占用严重超出硬件资源限制的问题，因此对应配置无法成功编译或执行。

因此，本文提出的优化 Pass 集合并非通过固定组合获得性能提升，而是通过将 Pass 开关纳入自动调优空间，使 DeepGen 能够根据算子形状、分块配置和目标硬件平台选择合适的优化组合。实验结果说明，不存在对所有场景都最优的固定 Pass 组合，自动调优机制对于发挥各优化 Pass 的条件性收益具有必要性。

% =============================================================
%                     4.5 计算精确度验证
% =============================================================
\section{计算精确度验证}\label{precision}

本节实验分别针对矩阵乘法算子和注意力机制进行测试。对于矩阵乘法，选取了三组方阵规模 $M=N=K$，分别为 512、1024 和 2048；对于注意力机制，则采用 $BHSD$ 形式描述输入规模，其中 $B$、$H$、$S$、$D$ 分别表示批大小、注意力头数、序列长度和头维度，测试配置包括 $(1,12,1024,64)$、$(1,16,2048,64)$ 和 $(1,32,4096,128)$。实验中均以 PyTorch 官方实现的计算结果作为参考基准，并采用相对 L2 范数误差衡量 DeepGen 输出与基准结果之间的差异。

\begin{table}[htbp]
\caption{矩阵乘法结果值的相对 L2 范数误差}\label{tab:table7}
\vspace{0.5em}\centering\wuhao
\renewcommand{\arraystretch}{1.35}
\begin{tabular}{ccc}
\toprule[1.5pt]
矩阵形状($M=N=K$) & NVIDIA & 海光 \\ \midrule[1pt]
512 & $4.457964\times 10^{-7}$ & $5.516395\times 10^{-7}$ \\
1024 & $6.003544\times 10^{-7}$ & $2.695841\times 10^{-6}$ \\ 
2048 & $8.012432\times 10^{-7}$ & $4.896582\times 10^{-6}$ \\ \bottomrule[1.5pt]
\end{tabular}
\vspace{\baselineskip}
\end{table}


\begin{table}[htbp]
\caption{多头注意力机制结果值的相对 L2 范数误差}\label{tab:table8}
\vspace{0.5em}\centering\wuhao
\renewcommand{\arraystretch}{1.35}
\begin{tabular}{ccc}
\toprule[1.5pt]
矩阵形状($BHSD$) & NVIDIA & 海光 \\ \midrule[1pt]
(1, 12, 1024, 64) & $3.869058\times 10^{-7}$ & $5.230882\times 10^{-7}$ \\
(1, 16, 2048, 64) & $6.390584\times 10^{-7}$ & $2.065086\times 10^{-6}$ \\ 
(1, 32, 4096, 128) & $8.330952\times 10^{-7}$ & $3.410693\times 10^{-6}$ \\
\bottomrule[1.5pt]
\end{tabular}
\vspace{\baselineskip}
\end{table}

由表 \ref{tab:table7} 可以看出，在矩阵乘法测试中，DeepGen 在 NVIDIA 平台上的相对 L2 范数误差分别为 $4.457964\times 10^{-7}$、$6.003544\times 10^{-7}$ 和 $8.012432\times 10^{-7}$，整体稳定在 $10^{-7}$ 量级；在海光平台上，误差分别为 $5.516395\times 10^{-7}$、$2.695841\times 10^{-6}$ 和 $4.896582\times 10^{-6}$，虽然随着矩阵规模增大略有上升，但总体仍保持在 $10^{-6}$ 量级以内。这说明 DeepGen 在矩阵乘法算子生成过程中，即使经历了分块、并行映射与访存优化等多种变换，其输出结果依然与 PyTorch 高度一致。

表 \ref{tab:table8} 给出了多头注意力机制的精度验证结果。可以看到，在 NVIDIA 平台上，三组配置下的误差分别为 $3.869058\times 10^{-7}$、$6.390584\times 10^{-7}$ 和 $8.330952\times 10^{-7}$，同样稳定在 $10^{-7}$ 量级；在海光平台上，误差分别为 $5.230882\times 10^{-7}$、$2.065086\times 10^{-6}$ 和 $3.410693\times 10^{-6}$，也始终处于较低水平。相较于矩阵乘法，多头注意力机制包含两次矩阵乘、缩放、Softmax 归一化以及输出重组等多个阶段，计算链路更长、数据流更复杂，因此其精度验证更能体现整个编译框架的数值可靠性。实验结果表明，DeepGen 在复杂融合算子场景下同样保持了良好的数值稳定性。

从整体趋势来看，NVIDIA 平台上的误差普遍低于海光平台，且基本稳定在 $10^{-7}$ 量级；海光平台随着问题规模增大，误差略有上升，但最大值仍远低于 $10^{-3}$ 的工业应用常用精度阈值。海光平台的误差虽然会随规模增长略有增加，但全尺寸下仍维持在较低范围内，说明该框架在不同硬件平台上均具备良好的数值一致性。造成平台间误差差异的主要原因，通常与底层指令实现、并行规约顺序以及不同硬件后端的浮点执行细节有关，但这类微小差异属于可接受的数值波动范围，并不会影响模型推理结果的有效性。

% =============================================================
%                4.6 端到端推理性能测试
% =============================================================
\section{端到端推理性能测试}

在本节中，本文将针 DeepGen 在端到端推理的效果方面进行实验评估，测量并记录模型在预填充（Prefill）阶段的单次推理延迟，并与 PyTorch 进行对比，以推理延迟计算加速比。实验平台分别选取了 NVIDIA A100/V100 GPU 和海光 K100 DCU，实验模型则选取了 BERT-Base、BERT-Large、GPT2-Small、GPT2-Medium 和 Llama2-7B 为测试对象，如表 \ref{tab:table9} 所示为各模型对于参数信息。

\begin{table}[htbp]
\caption{端到端实验模型参数}\label{tab:table9}
\vspace{0.5em}\centering\wuhao
\renewcommand{\arraystretch}{1.35}
\begin{tabular}{cccccc}
\toprule[1.5pt]
参数 & BERT-Base & BERT-Large & GPT2-Small & GPT2-Medium & Llama2-7B \\\midrule[1pt]
序列长度（SL） & 768 & 1024 & 1024 & 1024 & 2048 \\
注意力头维度（HD） & 64 & 64 & 64 & 64 & 128 \\
注意力头数（HN） & 12 & 16 & 12 & 16 & 32 \\
隐藏维度（Hidden dim） & 768 & 1024 & 768 & 1024 & 4096 \\
FNN 层维度 & 3072 & 4096 & 3072 & 4096 & 16384 \\
模型层数 & 12 & 16 & 12 & 24 & 32 \\
模型参数量 & 110M & 217M & 124M & 355M & 7B \\
\bottomrule[1.5pt]
\end{tabular}
\vspace{\baselineskip}
\end{table}

模型权重以随机方式初始化，其输入序列均采用随机生成的英文文本序列（词汇表大小为 32000）。其中对于 BERT-Base、BERT-Large、GPT2-Small、GPT2-Medium 模型的测试，性能主要以 PyTorch 为基线，并引入 TVM 作为在 NVIDIA 平台上的对照，而对于 Llama2-7B 模型的端到端测试主要在 NVIDIA A100 平台上进行，并且与 IREE、Tilelang、TVM 和 PyTorch 进行了对比。需要说明的是，本工作侧重于在普通核心如 CUDA 核心上的构建实现算子生成，因此对于 PyTorch 的评测未使用依赖于张量核的 FlashAttention 路径。

\begin{figure}[htbp]
\centering
\includegraphics[width=1\textwidth]{e2e_1}
\caption{BERT 和 GPT2 系列模型端到端推理实验}\label{e2e_1}
\vspace{\baselineskip}
\end{figure}

如图 \ref{e2e_1} 所示，BERT-Base、BERT-Large、GPT2-Small 和 GPT2-Medium 模型在不同平台上的测试结果表明，DeepGen 在多数配置下均能获得稳定的端到端性能收益。图中缺失的柱状结果表示对应配置下发生显存不足，无法完成评测，这也说明随着模型规模和批大小的增加，端到端推理对显存容量提出了更高要求。

对于 BERT 系列模型，DeepGen 在 A100、V100 和 K100 平台上均表现出较稳定的加速趋势。其中，BERT-Base 在小批量场景下收益更为明显，而随着批大小增大，加速比整体有所下降；BERT-Large 的变化趋势与 BERT-Base 基本一致，但由于模型规模更大，其可测试的批大小范围相对更小。这说明 DeepGen 对关键算子的优化能够有效转化为端到端收益，但随着批大小增大，运行时调度、内存访问和非关键算子开销在总延迟中的占比逐渐上升，从而削弱了局部优化在整体指标中的体现。

对于 GPT2-Small 和 GPT2-Medium，图 \ref{e2e_1} 同样显示 DeepGen 在多数配置下优于 PyTorch，并在 NVIDIA 平台上的多数配置下优于 TVM。总体来看，GPT2-Small 在 A100 和 V100 上保持了较稳定的正向加速，在 K100 上收益更加明显；GPT2-Medium 的趋势与之类似，但由于模型规模更大，其加速比相对更为收敛。这表明对于生成式 Transformer 模型而言，尽管完整推理过程包含更多层间依赖和状态传递，只要矩阵乘法和注意力机制等核心算子得到了有效优化，仍然能够在端到端层面获得可观收益。

\begin{figure}[htbp]
\centering
\includegraphics[width=0.7\textwidth]{e2e_2}
\caption{Llama 模型端到端推理实验}\label{e2e_2}
\vspace{\baselineskip}
\end{figure}

图\ref{e2e_2}展示了 Llama2-7B 在 A100 平台上的端到端推理结果。相较于 BERT 和 GPT2 系列，Llama2-7B 的模型规模更大、执行链路更长，因此其端到端加速比整体呈现更明显的收敛趋势。图中可以看出，DeepGen 相比 PyTorch 仍然保持了稳定的正向收益，并在多数配置下优于 TVM；与 IREE 相比，DeepGen 的整体表现更具优势，而 Tilelang 在该组实验中的性能则相对较低。这说明在大模型推理场景中，虽然 DeepGen 对多头注意力机制等关键算子具有较强优化能力，但最终端到端收益仍会受到层归一化、激活函数、张量重排以及运行时调度和内存管理开销的共同影响。

综合来看，DeepGen 在 BERT-Base、BERT-Large、GPT2-Small、GPT2-Medium 和 Llama2-7B 等模型的端到端推理中均取得了不同程度的性能提升。总体而言，DeepGen 在多数配置下优于 PyTorch，并在 NVIDIA 平台上的多数配置下优于 TVM；但随着模型规模和批大小增大，端到端加速比整体呈现一定收敛趋势。这说明端到端性能不仅取决于关键算子的局部优化，还受到系统级运行开销的共同影响。尽管如此，DeepGen 仍能在多平台、多模型场景下表现出稳定的正向收益，表明本文提出的基于 MLIR 的关键算子自动代码生成与优化方法具有较好的工程实用价值和跨平台适应能力。

% =============================================================
%                4.7 自动调优效率与搜索成本实验
% =============================================================
\section{自动调优效率与搜索成本实验}

为进一步评估 DeepGen 自动调优框架在实际部署过程中的可用性，本节从调优空间剪枝效果和搜索时间开销两个方面进行实验分析。实验分别在 NVIDIA A100 GPU 和海光 Z100 DCU 平台上进行，测试对象包括矩阵乘法算子和多头注意力机制算子。其中，矩阵乘法算子的输入规模统一设置为 $M=N=K=2048$，多头注意力机制算子的输入形状统一设置为 $BHSD=(1, 16, 2048, 64)$。为保证不同平台实验结果具有可比性，两个平台均采用相同的调优参数生成规则构造原始候选空间，并在此基础上结合各自硬件约束执行空间剪枝。

表 \ref{tab:table_1} 给出了 DeepGen 调优空间在剪枝前后的候选配置数量对比。原始候选空间仅依据算子计算结构和调优参数取值范围生成，主要保证候选参数组合在逻辑上具有基本合理性。例如，原始空间中已经考虑了线程束或波前中线程数量等平台相关基本参数，但尚未进一步约束线程块规模和内存资源等硬件限制，因此，原始候选空间规模较大。相比矩阵乘法，多头注意力机制算子包含更多可调超参数，因此其原始候选空间显著大于矩阵乘法算子。

\begin{table}[htbp]
\caption{调优空间剪枝前后候选配置数量对比}\label{tab:table_1}
\vspace{0.5em}\centering\wuhao
\renewcommand{\arraystretch}{1.35}
\begin{tabular}{ccccccc}
\toprule[1.5pt]
平台 & 算子类型 & 原始候选数 & 必要剪枝后 & 可选剪枝后 & 剪枝比例 & 最终候选数 \\ \midrule[1pt]
\multirow{2}{*}{NVIDIA A100} & MHA & 7,878,352 & 1,260,414 & 22,376 & 0.9971 & 22,376 \\
 & Matmul & 335,994 & 44,328 & 7,128 & 0.9787 & 7,128 \\ 
\multirow{2}{*}{海光 Z100} & MHA & 3,638,680 & 789,902 & 100,88 & 0.9972 & 10,088 \\
 & Matmul & 214,482 & 37,632 & 4,536 & 0.9788 & 4,536\\ \bottomrule[1.5pt]
\end{tabular}
\vspace{\baselineskip}
\end{table}

在生成原始候选空间后，DeepGen 进一步执行两级剪枝。首先进行可选剪枝，剔除不满足手动设置的限制条件，例如限制向量化拷贝下限和分块形状等。随后进行必要剪枝，根据线程数量、Bank 冲突预测以及缓存资源使用情况等条件进一步缩小搜索空间。由表 \ref{tab:table_1} 可以看出，剪枝策略能够显著降低候选配置数量。对于多头注意力机制算子，NVIDIA A100 平台上的候选数量由 7,878,352 个减少到 22,376 个，剪枝比例达到 99.71\%；海光 Z100 平台上的候选数量由 3,638,680 个减少到 10,088 个，剪枝比例达到 99.72\%。对于矩阵乘法算子，两个平台的剪枝比例也均达到约 97.88\%。这说明所设计的空间剪枝策略能够有效剔除大量无效或低潜力配置，从而避免对完整搜索空间进行暴力遍历。

在获得最终候选调优空间后，DeepGen 调优框架需要对候选配置进行代码生成、编译和性能测试，并从中选择性能最优的执行配置。本文在调优过程中采用多进程并行编译方式，以降低自动代码生成和编译阶段的时间开销。考虑到实验对比的一致性，两个平台仍统一设置为 80 个进程并行编译。

\begin{table}[htbp]
\caption{自动代码生成与调优时间开销}\label{tab:table_2}
\vspace{0.5em}\centering\wuhao
\renewcommand{\arraystretch}{1.35}
\begin{tabular}{cccccc}
\toprule[1.5pt]
平台 & 算子类型 & 最终候选数 & 编译时间（s） & 测试时间（s） & 总调优时间（s） \\ \midrule[1pt]
\multirow{2}{*}{NVIDIA A100} & MHA & 22,376 & 1381.80 & 1935.60 & 3492.12 \\
 & Matmul & 7,128 & 299.41 & 387.02 & 742.28 \\ 
\multirow{2}{*}{海光 Z100} & MHA & 10,088 & 3722.39 & 5755.18 & 9960.62 \\
 & Matmul & 4,536 & 1182.58 & 1674.63 & 2935.17 \\ \bottomrule[1.5pt]
\end{tabular}
\vspace{\baselineskip}
\end{table}

表 \ref{tab:table_2} 给出了 DeepGen 自动代码生成与调优时间开销统计结果。从表中可以看出，调优总时间与最终候选配置数量密切相关。多头注意力机制算子的最终候选空间明显大于矩阵乘法算子，因此其调优时间也更长。在 NVIDIA A100 平台上，多头注意力机制算子的最终候选数量为 22,376 个，总调优时间为 3492.12 秒；矩阵乘法算子的最终候选数量为 7,128 个，总调优时间为  742.28 秒。在海光 Z100 平台上，由于平台运行时系统以及测试执行效率等因素的影响，整体调优时间高于 NVIDIA A100 平台，其中多头注意力机制算子的总调优时间达到 9960.62 秒，矩阵乘法算子的总调优时间为 2935.17 秒。

进一步分析编译时间和测试时间可以发现，性能测试阶段在总调优时间中占据更高比例。这是因为测试阶段无法像编译阶段一样通过增加 CPU 并行进程无限制地降低开销。同时，为保证性能测量的稳定性，每个候选配置通常需要执行多次预热和重复测试，这也增加了整体搜索成本。相比之下，编译阶段可以通过增加编译进程数量的方式进一步优化。

总体来看，DeepGen 的空间剪枝策略能够显著压缩自动调优搜索空间，从而有效降低了暴力搜索带来的不可接受开销。然而，从最终候选数量和实际调优时间来看，即使经过多级剪枝，搜索空间仍然较大，完整遍历最终候选空间仍然需要较长时间。这表明当前调优框架虽然能够保证搜索的充分性和结果的稳定性，但在搜索效率方面仍存在进一步优化空间。



% =============================================================
%                4.8 本章小结
% =============================================================
\section{本章小结}

本章对 DeepGen 框架的实验效果进行了系统验证，主要从算子性能、数值精度和端到端推理性能三个方面展开。实验平台覆盖 NVIDIA 与海光 DCU 两类异构硬件，测试对象包括矩阵乘法、多头注意力以及 BERT、GPT2、Llama2 等典型 Transformer 模型。整体结果表明，DeepGen 在跨平台代码生成与性能优化方面具有较好的有效性。

在矩阵乘法实验中，DeepGen 在 NVIDIA A100 和海光 Z100 平台上均取得了与 cuBLAS、rocBLAS 接近甚至部分超过的性能表现，并在部分矩阵规模下实现反超。这说明，本文基于块级分层 IR 的算子建模方式，以及配套的降级信息推断、访存优化和自动调优机制，能够生成较好适配目标硬件特征的高性能实现。与 Triton、TVM 等工作的比较结果也表明，DeepGen 在保持统一中间表示和跨平台生成能力的同时，仍具备较强的性能竞争力。

在多头注意力实验中，DeepGen 在不同序列长度、头维度和批大小条件下均表现出较好的性能稳定性。在 A100 和 Z100 平台上，其在多数配置下取得了最优或接近最优的吞吐率，尤其在长序列和小头维度场景中优势更明显。这说明，DeepGen 对复杂 Transformer 核心算子的建模与优化方法是有效的。

在优化 Pass 消融实验中，本文验证了各优化策略对不同算子和硬件平台的影响存在显著差异，部分优化在提升性能的同时也可能引入额外开销。因此，DeepGen 将优化 Pass 的选择纳入自动调优空间，以根据具体场景适配最佳优化组合。

在数值精度方面，DeepGen 生成代码的输出结果与 PyTorch 官方实现保持高度一致。实验中，矩阵乘法和注意力算子的相对 L2 误差在 NVIDIA 平台上稳定在 $10^{-7}$ 量级，在海光平台上保持在 $10^{-6}$ 量级以内，说明相关优化并未破坏计算正确性。

在端到端推理测试中，DeepGen 在多模型、多平台场景下均取得了不同程度的性能提升，多数配置优于 PyTorch，并在 NVIDIA 平台上的多数配置下优于 TVM。综上，本文提出的基于 MLIR 的 Transformer 关键算子自动代码生成方法，能够在异构硬件上实现较高性能、较好精度和较强可移植性的代码生成。

在自动调优效率与搜索成本实验中，DeepGen 利用多种剪枝策略显著压缩了候选调优空间，减少了编译与测试开销。实验也表明，对于多头注意力机制等复杂算子，优化后的搜索空间仍然较大，自动调优时长仍有待进一步缩短。
